% GENERATED FILE. Edit canonical lessons, then run npm run generate:books.
% canonical-source-hash: fnv1a64:bb756dff53524720
% canonical-lessons: HI-C497-janm, HI-C497-garbhvati, HI-C497-prasav, HI-C497-navjat, HI-C497-paidahona

\chapter{A New Baby}
\label{ch:hi-a-new-baby}
\hlchaptermodality{497}

\begin{chapteropening}
\textbf{By the end of this chapter:} \emph{I can talk about the birth of a baby.}

Complete the last lesson of chapter 497: I can talk about the birth of a baby.
\end{chapteropening}

\section[\texorpdfstring{\hi{जन्म} (janm)}{जन्म (janm)}]{\hi{जन्म} (janm) — birth}

\label{lesson:HI-C497-janm}



\begin{quote}
Before the new one: say the Hindi for to mumble, then the Hindi for a bedroll.
\end{quote}

\subsection*{You'll want to know: \hi{जन्म}}
\textbf{\hi{जन्म}} — \emph{janm} — \textquotedblleft{}birth\textquotedblright{}.

\begin{quote}
\hi{मीरा} \hi{का} \hi{जन्म} \hi{एक} \hi{छोटे} \hi{गाँव} \hi{में} \hi{हुआ} \hi{था।}
\end{quote}

\begin{quote}
\emph{Mīrā kā janm ek choṭe gā̃v meṁ huā thā.}
\end{quote}

\begin{quote}
Mira was born in a small village.
\end{quote}

Masculine.

\begin{cousinweb}
\textbf{\hi{जन्म}} (\emph{janm}) is a tatsama, borrowed from Sanskrit \emph{janman}, \textquotedblleft{}birth, origin\textquotedblright{} (Monier-Williams). \emph{janmdin}, \textquotedblleft{}birthday\textquotedblright{}, is built on it.
\end{cousinweb}

\begin{grammarlens}[title={What you've built}]
One more word for describing what you lived and felt.
\end{grammarlens}

\subsection*{Your turn}
\begin{itemize}
\raggedright
  \item \emph{Say it:} \emph{janm}
  \item \emph{Say it:} \emph{janm}, once more
  \item \emph{Say it:} \emph{Mīrā kā janm ek choṭe gā̃v meṁ huā thā.} — Mira was born in a small village.
  \item \emph{Recall:} say the Hindi for to mumble, then the Hindi for a bedroll, then say \emph{janm} again
\end{itemize}

\begin{tcolorbox}[breakable,colback=teal!4,colframe=teal!35!black,title={Before you move on}]
What does \hi{जन्म} mean? (\textquotedblleft{}Birth\textquotedblright{}.) Say it once more.
\end{tcolorbox}
\section[\texorpdfstring{\hi{गर्भवती} (garbhvatī)}{गर्भवती (garbhvatī)}]{\hi{गर्भवती} (garbhvatī) — pregnant}

\label{lesson:HI-C497-garbhvati}



\begin{quote}
Before the new one: say the Hindi for a bedroll, then the Hindi for birth.
\end{quote}

\subsection*{You'll want to know: \hi{गर्भवती}}
\textbf{\hi{गर्भवती}} — \emph{garbhvatī} — \textquotedblleft{}pregnant\textquotedblright{}.

\begin{quote}
\hi{मेरी} \hi{भाभी} \hi{गर्भवती} \hi{हैं}, \hi{इसलिए} \hi{भारी} \hi{सामान} \hi{नहीं} \hi{उठातीं।}
\end{quote}

\begin{quote}
\emph{merī bhābhī garbhvatī haiṁ, isliye bhārī sāmān nahīṃ uṭhātīṃ.}
\end{quote}

\begin{quote}
My sister-in-law is pregnant, so she doesn't lift heavy things.
\end{quote}

\begin{cousinweb}
\textbf{\hi{गर्भवती}} (\emph{garbhvatī}) is a tatsama, borrowed from Sanskrit \emph{garbhavatī}, \textquotedblleft{}pregnant\textquotedblright{}, built on \emph{garbha}, \textquotedblleft{}womb\textquotedblright{}, and \emph{-vatī}, \textquotedblleft{}having\textquotedblright{} (Monier-Williams).
\end{cousinweb}

\begin{grammarlens}[title={What you've built}]
One more word for describing what you lived and felt.
\end{grammarlens}

\subsection*{Your turn}
\begin{itemize}
\raggedright
  \item \emph{Say it:} \emph{garbhvatī}
  \item \emph{Say it:} \emph{garbhvatī}, once more
  \item \emph{Say it:} \emph{merī bhābhī garbhvatī haiṁ, isliye bhārī sāmān nahīṃ uṭhātīṃ.} — My sister-in-law is pregnant, so she doesn't lift heavy things.
  \item \emph{Recall:} say the Hindi for a bedroll, then the Hindi for birth, then say \emph{garbhvatī} again
\end{itemize}

\begin{tcolorbox}[breakable,colback=teal!4,colframe=teal!35!black,title={Before you move on}]
What does \hi{गर्भवती} mean? (\textquotedblleft{}Pregnant\textquotedblright{}.) Say it once more.
\end{tcolorbox}
\section[\texorpdfstring{\hi{प्रसव} (prasav)}{प्रसव (prasav)}]{\hi{प्रसव} (prasav) — childbirth; delivery}

\label{lesson:HI-C497-prasav}



\begin{quote}
Before the new one: say the Hindi for birth, then the Hindi for pregnant.
\end{quote}

\subsection*{You'll want to know: \hi{प्रसव}}
\textbf{\hi{प्रसव}} — \emph{prasav} — \textquotedblleft{}childbirth; delivery\textquotedblright{}.

\begin{quote}
\hi{प्रसव} \hi{के} \hi{बाद} \hi{माँ} \hi{और} \hi{बच्चा} \hi{दोनों} \hi{ठीक} \hi{हैं।}
\end{quote}

\begin{quote}
\emph{prasav ke bād māṁ aur baccā donõ ṭhīk haiṁ.}
\end{quote}

\begin{quote}
After the delivery, mother and baby are both well.
\end{quote}

Masculine.

\begin{cousinweb}
\textbf{\hi{प्रसव}} (\emph{prasav}) is a tatsama, borrowed from Sanskrit \emph{prasava}, \textquotedblleft{}begetting, delivery, birth\textquotedblright{} (Monier-Williams). It is the word of hospitals and forms.
\end{cousinweb}

\begin{grammarlens}[title={What you've built}]
One more word for describing what you lived and felt.
\end{grammarlens}

\subsection*{Your turn}
\begin{itemize}
\raggedright
  \item \emph{Say it:} \emph{prasav}
  \item \emph{Say it:} \emph{prasav}, once more
  \item \emph{Say it:} \emph{prasav ke bād māṁ aur baccā donõ ṭhīk haiṁ.} — After the delivery, mother and baby are both well.
  \item \emph{Recall:} say the Hindi for birth, then the Hindi for pregnant, then say \emph{prasav} again
\end{itemize}

\begin{tcolorbox}[breakable,colback=teal!4,colframe=teal!35!black,title={Before you move on}]
What does \hi{प्रसव} mean? (\textquotedblleft{}Childbirth; delivery\textquotedblright{}.) Say it once more.
\end{tcolorbox}
\section[\texorpdfstring{\hi{नवजात} (navjāt)}{नवजात (navjāt)}]{\hi{नवजात} (navjāt) — newborn}

\label{lesson:HI-C497-navjat}



\begin{quote}
Before the new one: say the Hindi for pregnant, then the Hindi for childbirth.
\end{quote}

\subsection*{You'll want to know: \hi{नवजात}}
\textbf{\hi{नवजात}} — \emph{navjāt} — \textquotedblleft{}newborn\textquotedblright{}.

\begin{quote}
\hi{नवजात} \hi{बच्चा} \hi{दिन} \hi{भर} \hi{सोता} \hi{है।}
\end{quote}

\begin{quote}
\emph{navjāt baccā din bhar sotā hai.}
\end{quote}

\begin{quote}
A newborn baby sleeps all day long.
\end{quote}

\begin{cousinweb}
\textbf{\hi{नवजात}} (\emph{navjāt}) is a tatsama, borrowed from Sanskrit \emph{nava-jāta}, \textquotedblleft{}newly born\textquotedblright{}, built on \emph{nava}, \textquotedblleft{}new\textquotedblright{}, and \emph{jāta}, \textquotedblleft{}born\textquotedblright{} (Monier-Williams).
\end{cousinweb}

\begin{grammarlens}[title={What you've built}]
One more word for describing what you lived and felt.
\end{grammarlens}

\subsection*{Your turn}
\begin{itemize}
\raggedright
  \item \emph{Say it:} \emph{navjāt}
  \item \emph{Say it:} \emph{navjāt}, once more
  \item \emph{Say it:} \emph{navjāt baccā din bhar sotā hai.} — A newborn baby sleeps all day long.
  \item \emph{Recall:} say the Hindi for pregnant, then the Hindi for childbirth, then say \emph{navjāt} again
\end{itemize}

\begin{tcolorbox}[breakable,colback=teal!4,colframe=teal!35!black,title={Before you move on}]
What does \hi{नवजात} mean? (\textquotedblleft{}Newborn\textquotedblright{}.) Say it once more.
\end{tcolorbox}
\section[\texorpdfstring{\hi{पैदा} \hi{होना} (paidā honā)}{पैदा होना (paidā honā)}]{\hi{पैदा} \hi{होना} (paidā honā) — to be born}

\label{lesson:HI-C497-paidahona}



\begin{quote}
Before the new one: say the Hindi for childbirth, then the Hindi for newborn.
\end{quote}

\subsection*{You'll want to know: \hi{पैदा} \hi{होना}}
\textbf{\hi{पैदा} \hi{होना}} — \emph{paidā honā} — \textquotedblleft{}to be born\textquotedblright{}.

\begin{quote}
\hi{अरुण} \hi{की} \hi{बहन} \hi{के} \hi{घर} \hi{एक} \hi{बेटी} \hi{पैदा} \hi{हुई।}
\end{quote}

\begin{quote}
\emph{Aruṇ kī bahan ke ghar ek beṭī paidā huī.}
\end{quote}

\begin{quote}
A daughter was born to Arun's sister.
\end{quote}

Past (he …): \hi{वह} \hi{गाँव} \hi{में} \hi{पैदा} \hi{हुआ}, \emph{vah gā̃v meṁ paidā huā}.

\begin{cousinweb}
\textbf{\hi{पैदा} \hi{होना}} (\emph{paidā honā}) joins \emph{paidā}, usually traced to Persian \emph{paidā}, \textquotedblleft{}visible, born\textquotedblright{}, to \emph{honā}, \textquotedblleft{}to be\textquotedblright{}: to come into view. \emph{janm honā} says the same thing more formally.
\end{cousinweb}

\begin{grammarlens}[title={What you've built}]
One more word for describing what you lived and felt.
\end{grammarlens}

\subsection*{Your turn}
\begin{itemize}
\raggedright
  \item \emph{Say it:} \emph{paidā honā}
  \item \emph{Say it:} \emph{paidā honā}, once more
  \item \emph{Say it:} \emph{Aruṇ kī bahan ke ghar ek beṭī paidā huī.} — A daughter was born to Arun's sister.
  \item \emph{Recall:} say the Hindi for childbirth, then the Hindi for newborn, then say \emph{paidā honā} again
\end{itemize}

\begin{tcolorbox}[breakable,colback=teal!4,colframe=teal!35!black,title={Before you move on}]
What does \hi{पैदा} \hi{होना} mean? (\textquotedblleft{}To be born\textquotedblright{}.) Say it once more.
\end{tcolorbox}
